% Hierarchical partial-pooling structure
\begin{figure}[htbp]
  \centering
  \begin{tikzpicture}[x=1cm, y=1cm, every node/.style={font=\small}]
    % global hyperprior
    \node[hyper, fill=softgray] (global) at (5,3)
      {\,$\mu,\sigma$ \;\footnotesize(global hyperprior)\,};

    % segment-level parameters
    \node[latent] (uk) at (1.2,1.2) {UK};
    \node[latent] (de) at (3.7,1.2) {DE};
    \node[latent] (fr) at (6.3,1.2) {FR};
    \node[latent] (ot) at (8.8,1.2) {Other};

    % customers
    \node[box, minimum width=18mm] (cuk) at (1.2,-0.6) {\footnotesize many\\\footnotesize customers};
    \node[box, minimum width=16mm] (cde) at (3.7,-0.6) {\footnotesize few};
    \node[box, minimum width=16mm] (cfr) at (6.3,-0.6) {\footnotesize few};
    \node[box, minimum width=16mm] (cot) at (8.8,-0.6) {\footnotesize few};

    % edges: global -> segments (partial pooling)
    \draw[conn] (global) -- (uk);
    \draw[conn] (global) -- (de);
    \draw[conn] (global) -- (fr);
    \draw[conn] (global) -- (ot);

    % edges: segments -> customers
    \draw[conn] (uk) -- (cuk);
    \draw[conn] (de) -- (cde);
    \draw[conn] (fr) -- (cfr);
    \draw[conn] (ot) -- (cot);

    % shrinkage annotation
    \node[font=\footnotesize\itshape, accentorange, align=left, anchor=west] at (9.4,2.1)
      {\,shrinkage\\ \,toward global};
  \end{tikzpicture}
  \caption[Hierarchical partial pooling]{Hierarchical (partial-pooling) structure across country segments. Each segment's behavioural parameters are drawn from a shared global hyperprior, so estimates for data-sparse segments (DE, FR, Other) are pulled toward the global mean, borrowing strength from the data-rich UK segment. Complete pooling would collapse the middle layer; no pooling would remove the global hyperprior entirely.}
  \label{fig:hierarchy}
\end{figure}
